\section{Homogeneous polynomial maps and a posteriori enclosures}
\label{sec:enclosures}

The family above is one instance of a general counting reduction.  We state
the reusable part with the exact-decoding hypothesis exposed.  This avoids
suggesting that polynomial history counts automatically count distinct
words.

\begin{definition}[Disjoint polynomial count system]
Fix $\ell\ge2$.  A two-letter constant-length random-substitution level-count
system is \emph{disjoint polynomial} if its distinct inflation-word counts
satisfy
\begin{equation}
 A_{m+1}=P(A_m,B_m),\qquad B_{m+1}=Q(A_m,B_m),
 \label{eq:PQ-count}
\end{equation}
where $P,Q$ are nonzero homogeneous polynomials of degree $\ell$ with
nonnegative integer coefficients, and the identities have been justified by
disjoint level-block decoding.
\end{definition}

Put
\begin{equation}
 \phi(r)=\frac{Q(1,r)}{P(1,r)},\qquad
 \psi(r)=\log P(1,r).
 \label{eq:phi-psi}
\end{equation}
On any positive interval on which $P(1,r)>0$, both functions are continuous.

\begin{proposition}[Projective cocycle]
\label{prop:general-cocycle}
For a disjoint polynomial count system with $A_m,B_m>0$, let
$r_m=B_m/A_m$ and $G_m=\ell^{-m}\log A_m$.  Then
\begin{equation}
 r_{m+1}=\phi(r_m),\qquad
 G_{m+1}=G_m+\ell^{-(m+1)}\psi(r_m).
 \label{eq:general-cocycle}
\end{equation}
If the forward orbit lies in a compact interval on which $\psi$ is bounded,
then $G_m$ converges and
\begin{equation}
 G_\infty-G_M=\sum_{m=M}^{\infty}\ell^{-(m+1)}\psi(r_m).
 \label{eq:general-tail}
\end{equation}
Whenever an external entropy theorem identifies inflation-word entropy with
topological entropy, $G_\infty$ is that topological entropy.
\end{proposition}

\begin{proof}
Homogeneity gives
\[
 P(A_m,B_m)=A_m^\ell P(1,r_m),\qquad
 Q(A_m,B_m)=A_m^\ell Q(1,r_m).
\]
Division and logarithms yield \cref{eq:general-cocycle}.  Boundedness of
$\psi$ makes the discounted increment series absolutely convergent, and
telescoping gives \cref{eq:general-tail}.  The final identification is not a
counting consequence and is therefore stated conditionally.
\end{proof}

The next result turns rational interval arithmetic into a rigorous entropy
certificate.  It is deliberately formulated for arbitrary compact intervals;
monotonicity of $\phi$ or $\psi$ is useful computationally but is not assumed.

\begin{theorem}[Finite interval enclosure with invariant tail]
\label{thm:interval}
Fix integers $K\ge M$.  Suppose compact intervals
$I_M,I_{M+1},\ldots,I_K$ and a compact interval $J$ satisfy
\begin{equation}
 r_M\in I_M,\qquad
 \phi(I_m)\subseteq I_{m+1}\ (M\le m<K),\qquad
 I_K\subseteq J,\quad \phi(J)\subseteq J.
 \label{eq:interval-hyp}
\end{equation}
Assume $P(1,r)>0$ on their union.  Define
\begin{align}
 L_{M,K}&=\sum_{m=M}^{K-1}\ell^{-(m+1)}\min_{r\in I_m}\psi(r)
 +\frac{\ell^{-K}}{\ell-1}\min_{r\in J}\psi(r),\label{eq:enclosure-L}\\
 U_{M,K}&=\sum_{m=M}^{K-1}\ell^{-(m+1)}\max_{r\in I_m}\psi(r)
 +\frac{\ell^{-K}}{\ell-1}\max_{r\in J}\psi(r).
 \label{eq:enclosure-U}
\end{align}
Then
\begin{equation}
 L_{M,K}\le G_\infty-G_M\le U_{M,K}.
 \label{eq:enclosure-result}
\end{equation}
All endpoints are a posteriori quantities: replacing an interval by a
smaller enclosure, while re-verifying \cref{eq:interval-hyp}, can only sharpen
the corresponding extrema.
\end{theorem}

\begin{proof}
The inclusions in \cref{eq:interval-hyp} imply inductively that
$r_m\in I_m$ for $M\le m\le K$ and $r_m\in J$ for every $m\ge K$.
Bound the first $K-M$ terms of \cref{eq:general-tail} by the extrema on
$I_m$.  Bound every remaining term by the extrema on $J$ and use
\[
 \sum_{m=K}^{\infty}\ell^{-(m+1)}
 =\frac{\ell^{-K}}{\ell-1}.
\]
This proves both sides of \cref{eq:enclosure-result}.
\end{proof}

\begin{example}[Recovery of the family certificate]
For \cref{eq:substitution},
\[
 P(x,y)=x^\ell+y^\ell,\qquad Q(x,y)=x^py^q,\qquad
 \psi(r)=\log(1+r^\ell).
\]
The interval $J=[0,r_M]$ is forward invariant, and $\psi$ is increasing.
Taking $K=M$ in the evident invariant-tail specialization of
\cref{thm:interval} gives exactly \cref{eq:tail-certificate}; its lower
endpoint is zero, while strict positivity follows directly from the orbit.
\end{example}

\begin{remark}[What the general theorem does not say]
One must first prove \cref{eq:PQ-count} for \emph{sets of words}.  The
theorem neither detects collisions nor supplies recognizability.  It also
does not replace the hypotheses of the entropy-identification results in
\cite{Gohlke2020,Mitchell2024}.
\end{remark}
